\documentclass[12pt,a4paper]{report}

% ============================================================
% PACKAGES
% ============================================================
\usepackage[a4paper, top=1in, bottom=1in, left=1.5in, right=1in]{geometry}
\usepackage{setspace}
\usepackage[T1]{fontenc}
\usepackage[utf8]{inputenc}
\usepackage{mathptmx}
\usepackage{amsmath}
\usepackage{amssymb}
\usepackage{booktabs}
\usepackage{tabularx}
\usepackage{array}
\usepackage{graphicx}
\usepackage{caption}
\usepackage{subcaption}
\usepackage{fancyhdr}
\usepackage{titlesec}
\usepackage{tocloft}
\usepackage[hidelinks]{hyperref}
\usepackage{url}
\usepackage{enumitem}
\usepackage{parskip}
\usepackage{microtype}
\usepackage{xcolor}
\usepackage{listings}
\usepackage{float}
\usepackage{ragged2e}
\usepackage{tikz}
\usepackage{pgfplots}
\pgfplotsset{compat=1.18}
\usetikzlibrary{shapes.geometric, arrows.meta, positioning, fit,
                backgrounds, calc, decorations.pathreplacing, shadows,
                shapes.multipart, matrix}

% ============================================================
% COLOUR PALETTE (AWS-inspired for architecture diagrams)
% ============================================================
\definecolor{awsorange}{RGB}{255,153,0}
\definecolor{awsblue}{RGB}{35,103,178}
\definecolor{awsgreen}{RGB}{59,153,75}
\definecolor{awspurple}{RGB}{142,68,173}
\definecolor{awsgray}{RGB}{100,100,100}
\definecolor{awslightblue}{RGB}{213,232,255}
\definecolor{awslightgreen}{RGB}{213,245,218}
\definecolor{awslightorange}{RGB}{255,235,200}
\definecolor{awslightpurple}{RGB}{235,218,245}
\definecolor{sectionbg}{RGB}{245,245,245}
\definecolor{propblue}{RGB}{41,128,185}
\definecolor{propgreen}{RGB}{39,174,96}
\definecolor{propred}{RGB}{192,57,43}
\definecolor{propgray}{RGB}{149,165,166}

% ============================================================
% AIT FORMATTING
% ============================================================
\onehalfspacing
\setlength{\parindent}{0pt}
\setlength{\parskip}{8pt}

\titleformat{\chapter}[display]
  {\normalfont\fontsize{14}{16}\bfseries\centering}
  {\MakeUppercase{\chaptertitlename}\ \thechapter}{0pt}{\MakeUppercase}
\titlespacing*{\chapter}{0pt}{0pt}{20pt}

\titleformat{name=\chapter,numberless}[block]
  {\normalfont\fontsize{14}{16}\bfseries\centering}
  {}{0pt}{\MakeUppercase}
\titlespacing*{name=\chapter,numberless}{0pt}{0pt}{20pt}

\titleformat{\section}
  {\normalfont\fontsize{12}{14}\bfseries}{\thesection}{1em}{}
\titlespacing*{\section}{0pt}{14pt}{6pt}

\titleformat{\subsection}
  {\normalfont\fontsize{12}{14}\bfseries}{\thesubsection}{1em}{}
\titlespacing*{\subsection}{0pt}{12pt}{4pt}

\titleformat{\subsubsection}
  {\normalfont\fontsize{12}{14}\bfseries}{\thesubsubsection}{1em}{}
\titlespacing*{\subsubsection}{0pt}{10pt}{4pt}

\renewcommand{\cfttoctitlefont}{\normalfont\fontsize{14}{16}\bfseries}
\renewcommand{\contentsname}{CONTENTS}
\renewcommand{\listfigurename}{LIST OF FIGURES}
\renewcommand{\listtablename}{LIST OF TABLES}
\setlength{\cftbeforetoctitleskip}{0pt}
\setlength{\cftaftertoctitleskip}{14pt}
\renewcommand{\cftchapfont}{\bfseries}
\renewcommand{\cftchappagefont}{\bfseries}
\setlength{\cftbeforechapskip}{4pt}

\captionsetup[table]{position=above, labelsep=period,
  justification=raggedright, singlelinecheck=false, font=normalsize}
\captionsetup[figure]{position=below, labelsep=period,
  justification=justified, font=normalsize, singlelinecheck=false}

% Prevent overfull hbox in general text
\sloppy

\lstset{
  basicstyle=\ttfamily\small, breaklines=true, frame=single,
  backgroundcolor=\color{gray!8},
  xleftmargin=10pt, xrightmargin=10pt, aboveskip=8pt, belowskip=8pt
}

\graphicspath{{figures/}}

% ============================================================
% TIKZ STYLES
% ============================================================
\tikzset{
  % Pipeline boxes
  pbox/.style={rectangle, rounded corners=4pt, draw=#1!70, fill=#1!12,
               text centered, font=\small\bfseries, minimum height=0.65cm,
               minimum width=2.4cm, text width=2.2cm, drop shadow},
  % Arrow
  parrow/.style={-{Stealth[length=6pt]}, thick, color=awsgray},
  % Layer label
  layerlabel/.style={font=\footnotesize\bfseries, color=#1},
  % Component node
  compnode/.style={rectangle, rounded corners=3pt, draw=#1, fill=white,
                   font=\scriptsize, minimum height=0.55cm,
                   minimum width=1.8cm, text width=1.7cm, text centered},
  % Dashed bounding box
  boundbox/.style={rectangle, rounded corners=6pt, draw=#1!60, fill=#1!5,
                   dashed, thick},
  % Decision diamond
  decision/.style={diamond, draw=awsblue!70, fill=awslightblue,
                   font=\scriptsize, aspect=2.2, text centered,
                   minimum height=0.7cm, minimum width=2.0cm},
  % Evaluation row
  evalbox/.style={rectangle, rounded corners=3pt, draw=#1!70, fill=#1!10,
                  font=\small, minimum height=0.6cm, text width=3.5cm,
                  text centered},
}

% ============================================================
% DOCUMENT
% ============================================================
\begin{document}
\pagenumbering{roman}

% ============================================================
% TITLE PAGE
% ============================================================
\begin{titlepage}
\begin{center}
\vspace*{0.5cm}
{\fontsize{14}{18}\bfseries\MakeUppercase{%
TechPulse: A Real-Time Hybrid Retrieval-Augmented Generation\\
System for Emerging Technology Intelligence}\\[4pt]
\normalsize\mdseries\textit{Project Proposal}}

\vspace{1.8cm}
{\normalsize by}

    \vspace{1em}

    \begin{tabular}{cc}
        Aye Khin Khin Hpone (Yolanda Lim) & Dechathon Niamsa-Ard \\
        {\small st125970} & {\small st126235}
    \end{tabular}

\vfill

{\normalsize AT82.9002 Selected Topic: Data Engineering and MLOps}\\[6pt]
{\normalsize Instructor: Dr.\ Chantri Polprasert}\\[4pt]
{\normalsize Teaching Assistants: Swaraj Bhanja, Suryansh Srivastava}

\vfill
{\normalsize
Asian Institute of Technology\\[4pt]
School of Engineering and Technology\\[4pt]
Thailand\\[8pt]
February 2026
}
\end{center}
\end{titlepage}

% ============================================================
% ABSTRACT
% ============================================================
\chapter*{Abstract}
\addcontentsline{toc}{chapter}{ABSTRACT}

The rapid expansion of emerging technologies has created a fragmented knowledge ecosystem where authoritative scholarly research coexists with fast-moving industry discussions. Synthesizing accurate, timely, and verifiable intelligence across these heterogeneous sources presents a significant challenge for technology professionals. Large language models (LLMs) demonstrate advanced generative capabilities but are constrained by static training corpora, producing hallucinated or outdated outputs when queried about recent developments. Retrieval-Augmented Generation (RAG) offers a principled remedy by conditioning generation on externally retrieved evidence; however, baseline vector-only retrieval often suffers from contextual imprecision and retrieval noise.

This proposal introduces TechPulse, a fully AWS-native Hybrid Retrieval-Augmented Generation system designed to optimize grounded response generation under operational latency and cost constraints. The proposed system will integrate batch ingestion from ArXiv and near-real-time ingestion from Hacker News, deterministic preprocessing, token-based chunking, and semantic embedding using the MiniLM transformer. Embeddings will be indexed in Amazon Aurora PostgreSQL with pgvector, supporting integrated relational filtering and vector similarity search.

Retrieval will be structured as a multi-stage pipeline combining metadata-constrained filtering, cosine similarity search, and lightweight reranking incorporating semantic similarity, lexical overlap, and recency weighting. A controlled experiment will compare baseline vector-only retrieval against the proposed hybrid configuration under identical corpus and embedding conditions, evaluating context precision, faithfulness, p95 latency, and cost per query. The proposal outlines the system architecture, data engineering pipeline, MLOps deployment strategy, and evaluation plan. Specific implementation details remain subject to refinement as the project progresses.

\medskip
\noindent\textbf{Keywords:} Retrieval-Augmented Generation; Hybrid Information Retrieval; Vector Databases; Cloud-Native Architecture; MLOps; Semantic Embedding; Emerging Technology Intelligence.

% ============================================================
% FRONT MATTER
% ============================================================
\newpage\tableofcontents
\newpage\addcontentsline{toc}{chapter}{LIST OF TABLES}\listoftables
\newpage\addcontentsline{toc}{chapter}{LIST OF FIGURES}\listoffigures

\newpage
\pagenumbering{arabic}

% ============================================================
% CHAPTER 1: INTRODUCTION
% ============================================================
\chapter{Introduction}

\section{Problem Description and Motivation}
\label{sec:problem}

\textbf{Motivating scenario.} A machine learning engineer at a technology company needs to stay current with advances in transformer efficiency. She must manually monitor ArXiv preprints, scan Hacker News threads, and cross-reference multiple sources—a process that is time-consuming, fragmented, and prone to missing recent developments. An intelligent assistant capable of synthesizing both scholarly and community-sourced knowledge into a single, citation-grounded response would significantly accelerate her workflow. TechPulse is designed precisely for this use case.

The exponential growth of emerging technologies—including artificial intelligence, machine learning, distributed systems, and advanced computing—has produced an unprecedented volume of technical knowledge distributed across heterogeneous platforms. Scholarly repositories such as ArXiv provide structured, peer-reviewed or preprint research outputs contributing to theoretical and methodological advancement. Community-driven platforms and industry discussion forums such as Hacker News reflect near-real-time practitioner experiences and emerging tool ecosystems. However, the fragmentation of these knowledge regimes creates an information landscape that is increasingly difficult for technology professionals to synthesize effectively and in a timely manner.

Purely generative large language model (LLM) assistants are insufficient for emerging technology intelligence tasks due to their static knowledge cutoff, susceptibility to hallucinated responses, lack of verifiable evidence grounding, and inability to integrate with continuously updating data streams (Ji et al., 2023). Although Retrieval-Augmented Generation (RAG) architectures partially address these limitations by incorporating external context, baseline vector-only retrieval approaches may still retrieve semantically adjacent but contextually irrelevant documents. Such systems often fail to incorporate temporal relevance signals, structured metadata constraints, or recency-aware ranking mechanisms, which are essential for accurately capturing rapidly evolving technology trends.

Moreover, many existing RAG research prototypes focus primarily on retrieval accuracy without addressing production-readiness: deterministic pipeline design, idempotent ingestion, monitoring and observability, deployment automation, and cloud cost governance are frequently neglected. This limits their applicability in real-world operational environments.

The central engineering and research challenge addressed in this proposal is therefore:

\begin{quote}
\textbf{\textit{How can a scalable, continuously updating hybrid retrieval system be designed and deployed to synthesize emerging technology trends from heterogeneous scholarly and industry sources while satisfying latency and cost constraints?}}
\end{quote}

This project is motivated by three converging factors. First, technology professionals increasingly require timely and citation-grounded intelligence that synthesizes insights from both scholarly research and rapidly evolving industry discussions. Second, mature serverless cloud infrastructures now enable cost-effective deployment of production-grade data and machine learning pipelines. Third, a gap remains in the Retrieval-Augmented Generation (RAG) literature between retrieval-focused research prototypes and operationally robust systems capable of continuous ingestion, evaluation, and deployment in real-world environments.

\section{Research Questions}
\label{sec:rq}

To address the central problem, this study investigates the following research questions:

\begin{enumerate}[leftmargin=0.5in]
  \item \textbf{Data Engineering.} How can an AWS-native pipeline be designed to reliably ingest and preprocess heterogeneous scholarly and near-real-time industry data in a deterministic and idempotent manner?

  \item \textbf{Retrieval Optimization.} Does hybrid retrieval—combining metadata filtering, vector similarity search, and lightweight reranking—improve retrieval precision and grounding faithfulness compared with baseline vector-only retrieval?

  \item \textbf{Operational Trade-offs.} What trade-offs arise between retrieval precision, p95 latency, and cost per query under realistic workload conditions?

  \item \textbf{MLOps Integration.} How can deployment, monitoring, version control, and cost governance be integrated to ensure the reliability and maintainability of a production-grade RAG system on AWS?
\end{enumerate}

\section{Objectives of the Study}
\label{sec:objectives}

The primary objective is to design, implement, and evaluate a hybrid Retrieval-Augmented Generation system capable of synthesizing real-time emerging technology intelligence within a scalable, cost-aware cloud-native architecture. The specific objectives are:

\begin{enumerate}[leftmargin=0.5in]
  \item To construct an AWS-native data engineering pipeline for batch and near-real-time ingestion from ArXiv and Hacker News.
  \item To implement and compare baseline vector-only retrieval against hybrid retrieval with metadata filtering and reranking-lite optimization.
  \item To evaluate system performance across context precision, faithfulness, p95 latency, and cost per query.
  \item To deploy and monitor the system using serverless AWS infrastructure consistent with MLOps best practices (Sculley et al., 2015; Breck et al., 2017).
\end{enumerate}

\section{Significance of the Study}
\label{sec:significance}

This project contributes at the intersection of data engineering, natural language processing, and MLOps. From a technical perspective, it demonstrates how retrieval optimization can be integrated within a production-grade cloud architecture without relying on computationally expensive cross-encoder rerankers or specialized vector database infrastructure. The system provides a practical framework for grounded intelligence synthesis across academic and industry domains while maintaining architectural simplicity suitable for scalable deployment. From a practical perspective, the resulting system offers a blueprint for citation-supported technology intelligence synthesis that is operationally maintainable, cost-governed, and deployable within resource-constrained environments.


Chapter 1 introduces the problem description, motivation, research questions, and objectives of the study. Chapter 2 reviews relevant literature on Retrieval-Augmented Generation, hybrid retrieval strategies, vector databases, and MLOps practices. Chapter 3 presents the proposed methodology, technical approach, and technical innovation strategy. Chapter 4 describes the system architecture and data engineering pipeline design, including alignment with the AWS Well-Architected Framework. Chapter 5 outlines the experimental plan and evaluation metrics. Chapter 6 summarizes expected contributions and scope boundaries. References follow.

% ============================================================
% CHAPTER 2: BACKGROUND AND RELATED WORK
% ============================================================
\chapter{Background and Related Work}

\section{Retrieval-Augmented Generation}
\label{sec:lit_rag}

Retrieval-Augmented Generation (RAG) was introduced by Lewis et al. (2020) as a framework for augmenting sequence-to-sequence generation with a dense retrieval component, enabling the model to access external documents at inference time. This approach significantly improves performance on knowledge-intensive tasks by conditioning generation on retrieved evidence rather than relying solely on parametric memory. Despite their scale, LLMs remain prone to hallucination when responding to domain-specific or time-sensitive queries (Ji et al., 2023), making retrieval-based augmentation a practically important technique.

Recent surveys (Gao et al., 2023; Zhang et al., 2024) categorize RAG systems into four architectural types: (1) single-stage dense retrieval; (2) hybrid lexical--semantic retrieval; (3) multi-stage retrieval with reranking; and (4) iterative or self-reflective retrieval. While baseline dense retrieval improves factual grounding, empirical benchmarks such as BEIR (Thakur et al., 2021) demonstrate that semantic similarity alone is insufficient for high-precision retrieval in dynamic or heterogeneous corpora. Foundational retrieval-augmented systems such as Atlas (Izacard et al., 2022) and RETRO (Borgeaud et al., 2022) demonstrate retrieval augmentation at scale but require computationally intensive pretraining pipelines unsuited to cost-constrained production environments.

\section{Embedding Models and Vector Search}
\label{sec:lit_embed}

Dense retrieval relies on embedding models that transform text into high-dimensional vector representations. Sentence-transformer architectures such as MiniLM (Wang et al., 2020) provide computationally efficient embeddings with favorable trade-offs between semantic expressiveness and inference latency. PostgreSQL with the pgvector extension enables cosine similarity search alongside SQL-based metadata filtering within a single relational storage layer, reducing architectural complexity compared to dedicated vector databases such as FAISS or Milvus. This integrated approach is increasingly favored for production RAG deployments (Gao et al., 2023).

\section{Hybrid Retrieval and Reranking}
\label{sec:lit_hybrid}

Hybrid retrieval combines lexical matching (e.g., BM25 or keyword filtering) with semantic similarity search to improve retrieval precision. Thakur et al. (2021) demonstrate that hybrid approaches consistently outperform purely dense retrieval across heterogeneous benchmarks. Reranking further refines candidate relevance using additional contextual signals. Cross-encoder rerankers achieve high accuracy but incur substantial computational overhead (Nogueira \& Cho, 2019); lightweight multi-factor reranking—combining similarity score, recency weighting, and keyword overlap—offers comparable grounding improvements at lower latency (Gao et al., 2023). In dynamic knowledge domains, recency-aware retrieval is particularly important to avoid outdated context injection.

\section{RAG Evaluation Methodologies}
\label{sec:lit_eval}

Evaluating RAG systems requires assessing both retrieval quality and generation faithfulness simultaneously. Es et al. (2023) introduced RAGAS, a reference-free evaluation framework that measures faithfulness, answer relevance, context precision, and context recall using automated LLM-based scoring, enabling scalable evaluation without large human annotation pools. Complementary work by Saad-Falcon et al. (2023) proposed ARES, which uses a small set of human-annotated examples to train lightweight judges capable of scoring RAG systems across multiple quality dimensions. These frameworks demonstrate that even compact evaluation sets—in the range of 20 to 50 annotated examples—can yield statistically meaningful comparisons between retrieval configurations when evaluation dimensions are clearly scoped and metrics are carefully defined (Es et al., 2023). This finding provides empirical support for the evaluation scale adopted in TechPulse.

\section{Data and Concept Drift in Retrieval Systems}
\label{sec:lit_drift}

In dynamic knowledge domains, retrieval system performance can degrade over time due to data drift—a shift in the statistical properties of incoming documents relative to the indexed corpus—or concept drift, where the semantic meaning of user queries evolves independently of corpus updates (Gama et al., 2014). Lu et al. (2018) provide a comprehensive taxonomy of drift detection methods, categorizing approaches into statistical process control, windowed distribution comparison, and ensemble-based disagreement detection. In retrieval-augmented systems specifically, drift can manifest as declining context precision when newly ingested documents diverge from the embedding space learned during the original indexing phase. Addressing drift requires both monitoring mechanisms that can detect distributional shifts in retrieval scores over time, and remediation actions such as corpus re-indexing, embedding model refresh, or query routing adjustments (Bayram et al., 2022). TechPulse incorporates a lightweight embedding distribution monitoring strategy informed by these findings.

\section{Hallucination Detection and Faithfulness Verification}
\label{sec:lit_hallucination}

Despite retrieval augmentation, generative models can still produce hallucinated content by inferring beyond the supplied context, over-generalizing retrieved passages, or introducing spurious cross-document associations (Ji et al., 2023). Manakul et al. (2023) proposed SelfCheckGPT, a consistency-based approach that detects hallucinations by sampling multiple outputs and measuring inter-sample consistency—a technique applicable without access to ground-truth references. Min et al. (2023) proposed FActScoring, which decomposes generated responses into atomic claims and individually verifies each claim against a reference corpus, offering fine-grained faithfulness assessment. At a system design level, Shi et al. (2023) demonstrate that irrelevant context injected alongside relevant passages can distract generative models and induce hallucinated responses even when correct information is present—motivating strict top-$k$ selection and high-precision retrieval as the primary anti-hallucination mechanism in TechPulse.

\section{MLOps and Production Machine Learning Systems}
\label{sec:lit_mlops}

Sculley et al. (2015) identify hidden technical debt as a central challenge in ML systems, emphasizing monitoring, reproducibility, and lifecycle management. MLOps extends DevOps principles to ML workflows, incorporating continuous integration, deployment automation, model versioning, and observability (Breck et al., 2017). Cloud-native, serverless architectures enable scalable and maintainable deployment with cost governance (Zhang et al., 2020; AWS, 2023). While RAG research frequently focuses on retrieval effectiveness, production-readiness—including pipeline idempotency, schema versioning, and budget-aware inference—remains underexplored.

\section{Research Gap}
\label{sec:gap}

The reviewed literature establishes six key insights: (1) RAG improves factual grounding but baseline dense retrieval may lack precision; (2) hybrid retrieval outperforms vector-only search in heterogeneous settings; (3) lightweight reranking can improve relevance without prohibitive overhead; (4) production ML systems require robust MLOps practices beyond model accuracy; (5) even compact evaluation sets of 20--50 annotated examples can yield statistically meaningful RAG comparisons when evaluation dimensions are clearly scoped; and (6) hallucination in RAG systems can persist despite retrieval augmentation, necessitating explicit faithfulness verification mechanisms. However, no existing work integrates all six dimensions into a unified, cost-aware, cloud-native system explicitly targeting dynamic emerging technology intelligence. This project addresses that gap.

Table~\ref{tab:related_work} summarizes representative related work and positions TechPulse relative to prior systems.

\begin{table}[H]
\caption{Comparison of Related Work and Positioning of TechPulse.}
\label{tab:related_work}
\centering
\begin{tabular}{p{2.8cm}p{2.2cm}p{1.6cm}p{2.0cm}p{1.7cm}p{1.5cm}}
\toprule
\textbf{System} & \textbf{Retrieval} & \textbf{Rerank} & \textbf{Real-Time} & \textbf{Cost Gov.} & \textbf{MLOps} \\
\midrule
RAG (Lewis et al., 2020) & Dense & No & No & No & No \\
Atlas (Izacard et al., 2022) & Dense & No & No & No & No \\
RETRO (Borgeaud et al., 2022) & Dense & No & No & No & No \\
Typical RAG Prototype & Dense & Optional & Rare & No & Minimal \\
\textbf{TechPulse (proposed)} & \textbf{Hybrid} & \textbf{Yes} & \textbf{Yes} & \textbf{Yes} & \textbf{Full} \\
\bottomrule
\end{tabular}
\end{table}

\section{Chapter Summary}

This chapter reviewed the evolution of RAG architectures, embedding-based retrieval, hybrid retrieval optimization, MLOps principles, RAG evaluation methodologies, drift detection strategies, and hallucination verification approaches. The literature confirms that hybrid retrieval with lightweight reranking and production-grade MLOps integration represents an underexplored but practically important direction, and that compact but well-scoped evaluation sets can yield statistically meaningful RAG comparisons. The next chapter presents the methodology proposed to address this gap.

% ============================================================
% CHAPTER 3: PROPOSED METHODOLOGY
% ============================================================
\chapter{Proposed Methodology}

\section{Overview of Technical Approach}
\label{sec:method_overview}

TechPulse is proposed as a fully serverless, AWS-native Hybrid Retrieval-Augmented Generation system designed to synthesize emerging technology intelligence from heterogeneous sources. The technical approach is organized around five interconnected components: (1) continuous heterogeneous data ingestion; (2) deterministic preprocessing and semantic embedding; (3) hybrid multi-stage retrieval; (4) context-bound grounded generation; and (5) production-grade MLOps deployment and monitoring.

The core hypothesis is that hybrid retrieval—combining structured metadata filtering, dense cosine similarity search, and lightweight multi-factor reranking—will improve contextual precision and grounding faithfulness compared to baseline vector-only retrieval, while maintaining acceptable p95 latency and remaining within defined cost thresholds. This hypothesis will be tested through a controlled experiment comparing both configurations under identical conditions.

It is important to note that this proposal represents a preliminary plan. Specific implementation choices—including embedding model selection, chunking parameters, reranking weight tuning, and infrastructure configuration—may be refined as the project progresses and empirical results become available.

\section{Why This Approach is Suitable}
\label{sec:method_rationale}

Three considerations motivate the hybrid RAG approach for this specific problem domain.

\textbf{Domain characteristics.} Emerging technology knowledge is temporally dynamic and distributed across heterogeneous sources. Vector-only retrieval treats all documents equally regardless of recency or source authority; hybrid retrieval with recency-aware reranking explicitly addresses temporal drift, which is critical in rapidly evolving AI/ML domains where a paper from 18 months ago may already be superseded.

\textbf{Operational constraints.} The system must operate within cost and latency bounds suitable for a research prototype. The lightweight linear reranking function (Equation~\ref{eq:reranking}) adds $O(k)$ computation over the top-$N$ filtered candidates—negligible compared to the $O(Nd)$ baseline search—and avoids the latency penalties of cross-encoder rerankers, which typically require a forward pass through a BERT-class model per candidate.

\textbf{Production feasibility.} A fully serverless AWS architecture—Lambda, Aurora PostgreSQL, API Gateway, EventBridge—eliminates the need to manage dedicated compute clusters or self-hosted vector databases. This enables a student team to deploy, iterate, and monitor a production-grade system within the course timeline and a prototype budget of under USD~15/month.

\section{Technical Innovation and Optimization Strategy}
\label{sec:innovation}

The technical novelty of TechPulse lies not in any single component but in their integration into a coherent, cost-governed MLOps system. Specifically, four aspects represent the technical contribution of this project.

\textbf{Recency-aware hybrid reranking for dynamic knowledge domains.} The weighted linear reranking function (Equation~\ref{eq:reranking}) fuses semantic similarity, lexical keyword overlap, and exponential temporal decay into a single interpretable score. Unlike cross-encoder rerankers that are domain-agnostic and computationally expensive, this formulation is specifically designed for domains where content freshness is a first-class retrieval signal. The exponential decay parameter $\lambda$ provides a principled, tunable mechanism for controlling how rapidly older documents lose retrieval priority.

\textbf{Unified vector--relational retrieval without dedicated infrastructure.} By extending Amazon Aurora PostgreSQL with pgvector, TechPulse achieves hybrid retrieval (structured SQL metadata filtering + cosine vector search) within a single managed database, eliminating the operational complexity of running a separate vector store. This design choice is an explicit optimization: it reduces architectural surface area, simplifies deployment automation, and enables cost-governed scaling through Aurora Serverless.

\textbf{Idempotent ingestion pipeline with SHA-256 content-addressed deduplication.} The SHA-256 hashing strategy applied to canonicalized document fields ensures that the pipeline can be re-run safely at any time without creating duplicate records. This is an often-overlooked but practically important data engineering innovation—it enables safe experiment replay, crash recovery, and corpus versioning without requiring a separate deduplication service.

\textbf{Budget-aware inference with programmatic halt.} The integration of AWS Budgets alerts into the inference path—where invocations are halted programmatically when monthly cost thresholds are approached—represents an operational optimization strategy not typically addressed in RAG research prototypes. This transforms cost from a passive monitoring concern into an active constraint enforced at runtime, consistent with production-grade MLOps principles.

\section{Baseline Retrieval Configuration}
\label{sec:baseline}

The baseline configuration performs dense cosine similarity search across the full indexed corpus. Given a query embedding $\mathbf{e}_q$ and document embeddings $\mathbf{e}_i$, cosine similarity is computed as:

\begin{equation}
  \text{Sim}(q, i) = \frac{\mathbf{e}_q \cdot \mathbf{e}_i}{\|\mathbf{e}_q\|\|\mathbf{e}_i\|}
  \label{eq:cosine}
\end{equation}

Top-$N$ documents are selected purely on similarity ranking. Under exact search, computational complexity scales as $O(Nd)$, where $N$ is the number of indexed chunks and $d$ the embedding dimension. This configuration serves as the experimental control for retrieval performance comparison.

\section{Proposed Hybrid Retrieval Configuration}
\label{sec:hybrid}

The hybrid configuration introduces three sequential refinement stages applied before generation.

\textbf{Stage 1 — Metadata Filtering.} Structured SQL constraints reduce the candidate space based on document recency (a 180-day window is proposed as an initial parameter, subject to tuning), source type, and keyword presence. This pre-filtering stage reduces retrieval noise prior to semantic ranking. The 180-day window is intended to balance recency with coverage of foundational work; this threshold will be evaluated empirically.

\textbf{Stage 2 — Vector Similarity Search.} Cosine similarity (Equation~\ref{eq:cosine}) is applied over the filtered candidate subset rather than the full corpus, reducing both computational overhead and retrieval noise.

\textbf{Stage 3 — Reranking-Lite.} A lightweight weighted scoring function refines candidate ordering:

\begin{equation}
  \text{Score}_i = \alpha \cdot \text{Sim}_i + \beta \cdot \text{KeywordOverlap}_i + \gamma \cdot \text{RecencyWeight}_i
  \label{eq:reranking}
\end{equation}

where $\alpha + \beta + \gamma = 1$. Initial weights are proposed as $\alpha = 0.6$, $\beta = 0.2$, $\gamma = 0.2$. The rationale for this initialization is threefold. First, dense semantic similarity ($\alpha$) is assigned the largest weight because it captures latent semantic relatedness that neither keyword overlap nor recency can replicate; this is consistent with the empirical finding across BEIR benchmarks (Thakur et al., 2021) that dense retrieval dominates for open-domain tasks. Second, lexical keyword overlap ($\beta$) is assigned a secondary weight because technology queries frequently contain specialized terminology (model names, acronyms, version identifiers) where exact-match signals provide complementary precision gains not captured by embedding-space proximity. Third, recency ($\gamma$) is assigned the same secondary weight as keyword overlap, reflecting the temporal sensitivity of emerging technology intelligence without allowing recency to override semantic relevance entirely—a risk if $\gamma$ were set too high.

To verify that these initialized weights are appropriate and to systematically tune them, the following strategy will be applied during the evaluation phase. Grid search will be performed over the discrete weight space $\{0.4, 0.5, 0.6, 0.7\}$ for $\alpha$, with $\beta$ and $\gamma$ receiving equal shares of the remainder, spanning configurations such as (0.6, 0.2, 0.2), (0.7, 0.15, 0.15), and (0.5, 0.25, 0.25). For each configuration, context precision will be measured across the full 20--30 query evaluation set. The configuration yielding the highest mean precision without violating the p95 latency threshold will be selected as the final production weight configuration. This approach ensures that weights are empirically validated rather than assumed, while keeping the tuning procedure tractable within the project timeline.

Recency is modeled using exponential decay:

\begin{equation}
  \text{RecencyWeight}_i = e^{-\lambda \Delta t_i}
  \label{eq:recency}
\end{equation}

where $\Delta t_i$ is document age in days and $\lambda = 0.01$ (subject to tuning). This formulation ensures smooth temporal decay without abruptly excluding older but potentially relevant documents.

Figure~\ref{fig:retrieval_comparison} illustrates the conceptual difference between the baseline and hybrid retrieval pathways.

% ---- TikZ: Baseline vs Hybrid comparison diagram ----
\begin{figure}[H]
\centering
\begin{tikzpicture}[node distance=0.45cm and 1.4cm]

% ---- LEFT: BASELINE ----
\node[pbox=awsgray, minimum width=2.6cm] (bq) {User Query};
\node[pbox=awsblue, below=of bq, minimum width=2.6cm] (bvs) {Vector Similarity\\Search (Full Corpus)};
\node[pbox=awsgreen, below=of bvs, minimum width=2.6cm] (btopk) {Top-$k$ Selection};
\node[pbox=awsorange, below=of btopk, minimum width=2.6cm] (bgen) {Bedrock\\Generation};

\draw[parrow] (bq) -- (bvs);
\draw[parrow] (bvs) -- (btopk);
\draw[parrow] (btopk) -- (bgen);

\node[layerlabel=awsgray, above=0.3cm of bq, font=\small\bfseries] {\textbf{Baseline (Vector-Only)}};

\begin{scope}[on background layer]
\node[boundbox=awsgray, fit=(bq)(bvs)(btopk)(bgen), inner sep=8pt] {};
\end{scope}

% ---- RIGHT: HYBRID ----
\node[pbox=awsgray, minimum width=2.6cm, right=3.6cm of bq] (hq) {User Query};
\node[pbox=propred, below=of hq, minimum width=2.6cm] (hmf) {Metadata Filter\\(recency + source)};
\node[pbox=awsblue, below=of hmf, minimum width=2.6cm] (hvs) {Vector Similarity\\Search (Filtered)};
\node[pbox=awspurple, below=of hvs, minimum width=2.6cm] (hrr) {Reranking-Lite\\($\alpha$·Sim + $\beta$·Kw + $\gamma$·Rec)};
\node[pbox=awsgreen, below=of hrr, minimum width=2.6cm] (htopk) {Top-$k$ Selection};
\node[pbox=awsorange, below=of htopk, minimum width=2.6cm] (hgen) {Bedrock\\Generation};

\draw[parrow] (hq) -- (hmf);
\draw[parrow] (hmf) -- (hvs);
\draw[parrow] (hvs) -- (hrr);
\draw[parrow] (hrr) -- (htopk);
\draw[parrow] (htopk) -- (hgen);

\node[layerlabel=propblue, above=0.3cm of hq, font=\small\bfseries] {\textbf{Proposed Hybrid}};

\begin{scope}[on background layer]
\node[boundbox=propblue, fit=(hq)(hmf)(hvs)(hrr)(htopk)(hgen), inner sep=8pt] {};
\end{scope}

% annotation arrows
\draw[-{Stealth}, dashed, color=propred!70, thick]
  ([xshift=1.4cm]hmf.east) -- ++(0.5,0)
  node[right, font=\scriptsize, color=propred, text width=2.0cm]
  {Reduces candidate\\pool before search};

\draw[-{Stealth}, dashed, color=awspurple!70, thick]
  ([xshift=1.4cm]hrr.east) -- ++(0.5,0)
  node[right, font=\scriptsize, color=awspurple, text width=2.0cm]
  {Multi-factor\\relevance scoring};

\end{tikzpicture}
\caption{Comparison of Baseline Vector-Only Retrieval (left) and Proposed Hybrid Retrieval Pipeline (right). The hybrid pipeline adds a metadata filtering stage to reduce the candidate pool and a reranking-lite stage for multi-factor relevance scoring prior to final candidate selection.}
\label{fig:retrieval_comparison}
\end{figure}

\section{Prompt Construction, Grounded Generation, and Hallucination Verification}
\label{sec:prompt}

Retrieved top-$k$ documents will be concatenated into a structured prompt template that instructs the generative model to produce responses exclusively grounded in the supplied context. The prompt will explicitly prohibit unsupported speculation and require the model to acknowledge insufficient evidence when appropriate. Context length will be truncated to remain within model token limits, ensuring deterministic inference behavior. This approach mitigates hallucination risk and reinforces citation traceability. The exact prompt template and context truncation strategy will be finalized during implementation.

\textbf{Hallucination verification strategy.} Prompt-level constraints are a necessary but insufficient safeguard against hallucination. TechPulse will apply a three-layer verification approach to detect and flag potentially hallucinated outputs. First, at the prompt layer, the system instruction enforces strict evidence grounding and abstention on insufficient context (see Listing~\ref{lst:prompt} in Chapter~\ref{sec:prompt_ch4}). Second, at the evaluation layer, RAGAS faithfulness scoring will be applied to a sample of system outputs by decomposing each generated response into atomic claims and checking whether each claim is supported by at least one retrieved document. Any response where RAGAS faithfulness falls below a threshold of 0.8 will be flagged for manual review. Third, at the output layer, citation traceability will be enforced by requiring the model to tag each generated statement with a document reference identifier from the context block; responses that contain unanchored claims—statements with no corresponding source citation—will be intercepted and replaced with an abstention message. Together, these three layers provide defense-in-depth against hallucinated outputs, combining structural prompt design, automated scoring, and post-generation citation validation.

\section{Optimization Objective}
\label{sec:optimization}

The retrieval task is framed as a constrained optimization problem:

\begin{equation}
  \max \; \text{Precision}(q) \quad \text{subject to} \quad \text{Latency} \leq \tau
  \label{eq:opt}
\end{equation}

where $\tau$ is an initial p95 latency target of 2 seconds. This formulation makes the precision--latency trade-off explicit and provides a principled basis for tuning reranking depth and metadata filter selectivity.

\section{Tool Selection and Feasibility}
\label{sec:tools}

Table~\ref{tab:tools} summarizes the proposed tools and their justification. The full stack is serverless and AWS-managed, eliminating the need for dedicated cluster management and enabling cost-governed scaling within the project timeline.

\begin{table}[H]
\caption{Proposed Tools, Services, and Justification.}
\label{tab:tools}
\centering
\begin{tabular}{p{3.2cm}p{3.2cm}p{6.6cm}}
\toprule
\textbf{Component} & \textbf{Tool / Service} & \textbf{Justification} \\
\midrule
Ingestion orchestration & AWS Lambda + EventBridge & Serverless, scheduled/event-driven, no cluster management \\
Message queue & Amazon SQS + DLQ & Decouples ingestion from embedding; failure resilience \\
Raw storage & Amazon S3 (medallion layout) & Cost-effective, durable, prefix-based lineage tracking \\
Embedding model & MiniLM (sentence-transformers) & Low latency, strong semantic expressiveness, serverless-compatible \\
Vector + relational store & Aurora PostgreSQL + pgvector & Unified hybrid retrieval without dedicated vector DB overhead \\
Generative model & Amazon Bedrock (Claude Haiku) & Low-latency inference, managed, Asia-Pacific region availability \\
Frontend & React (S3 static hosting) & Lightweight, decoupled from backend \\
API routing & Amazon API Gateway & Managed, HTTPS, rate limiting \\
CI/CD & GitHub Actions + SAM/CDK & Reproducible IaC deployments, automated testing \\
Monitoring & CloudWatch + AWS Budgets + SNS & Observability, cost alerts, health checks \\
Evaluation & RAGAS + manual annotation & Standardized RAG quality metrics + qualitative grounding check \\
\bottomrule
\end{tabular}
\end{table}

Regarding project feasibility: the complete stack uses managed AWS services with free-tier or low-cost tiers available for prototype-scale usage. Estimated monthly operational cost under moderate evaluation workload (approximately 900 queries/month) is projected at USD~10--15, based on Bedrock token pricing, Lambda execution duration, and Aurora Serverless usage. This estimate will be validated during the implementation phase.

% ============================================================
% CHAPTER 4: SYSTEM ARCHITECTURE AND PIPELINE DESIGN
% ============================================================
\chapter{System Architecture and Pipeline Design}

\section{Architectural Overview}
\label{sec:arch_overview}

TechPulse is proposed as a six-layer, cloud-native system. Each layer is logically independent with well-defined interfaces, enabling modular development, testing, and future extension. Figure~\ref{fig:aws_arch} presents the full AWS service-level architecture, illustrating how the four functional zones—Sources, Ingestion, Processing, Serving, CI/CD, and Observability—interact using native AWS services with real AWS console icons. Figure~\ref{fig:six_layers} then presents the corresponding six-layer conceptual abstraction.

% ---- AWS Architecture Diagram (with official AWS service icons) ----
\begin{figure}[H]
\centering
\includegraphics[width=\textwidth]{aws_techpulse.png}
\caption{Full AWS Service-Level Architecture of TechPulse. The diagram shows the five functional zones: \textbf{Sources} (ArXiv API, Hacker News API); \textbf{Ingestion} (EventBridge scheduler triggering Ingestion Lambda); \textbf{Processing} (SQS Queue, S3 Raw Buckets, Preprocess + Embed Lambda, S3 Processed, Dead Letter Queue, Aurora + pgvector); \textbf{Serving} (S3 Static Hosting, React UI, API Gateway, RAG Orchestrator Lambda, Amazon Bedrock); \textbf{CI/CD} (GitHub Repo, GitHub Actions, SAM/CDK); and \textbf{Observability} (Health Check Lambda, CloudWatch, AWS Budgets, SNS). Solid arrows represent data and deployment flow; bidirectional arrows indicate API communication; dashed arrows represent monitoring and metrics signals.}
\label{fig:aws_arch}
\end{figure}

% ---- TikZ: Six-layer architectural stack ----
\begin{figure}[H]
\centering
\begin{tikzpicture}[node distance=0.0cm]

\def\lw{13.5cm}
\def\lh{0.90cm}

\node[rectangle, draw=awsgray!50, fill=awslightblue, minimum width=\lw,
      minimum height=\lh, font=\small\bfseries, text centered,
      rounded corners=2pt]
  (l6) {Layer 6 — Deployment, Monitoring \& Interaction \quad \textcolor{awsgray}{\scriptsize CloudWatch, Budgets, API GW, React}};

\node[rectangle, draw=awsgray!50, fill=awslightorange, minimum width=\lw,
      minimum height=\lh, font=\small\bfseries, text centered,
      rounded corners=2pt, below=0pt of l6]
  (l5) {Layer 5 — Generative Synthesis (Amazon Bedrock) \quad \textcolor{awsgray}{\scriptsize Claude Haiku, grounded prompting}};

\node[rectangle, draw=awsgray!50, fill=awslightpurple, minimum width=\lw,
      minimum height=\lh, font=\small\bfseries, text centered,
      rounded corners=2pt, below=0pt of l5]
  (l4) {Layer 4 — Hybrid Retrieval \& Reranking \quad \textcolor{awsgray}{\scriptsize 3-stage: filter $\to$ similarity $\to$ rerank}};

\node[rectangle, draw=awsgray!50, fill=awslightgreen, minimum width=\lw,
      minimum height=\lh, font=\small\bfseries, text centered,
      rounded corners=2pt, below=0pt of l4]
  (l3) {Layer 3 — Vector Storage \& Indexing (Aurora + pgvector) \quad \textcolor{awsgray}{\scriptsize cosine search + SQL filter}};

\node[rectangle, draw=awsgray!50, fill=awslightorange, minimum width=\lw,
      minimum height=\lh, font=\small\bfseries, text centered,
      rounded corners=2pt, below=0pt of l3]
  (l2) {Layer 2 — Preprocessing \& Feature Engineering \quad \textcolor{awsgray}{\scriptsize Normalize, chunk, SHA-256 hash}};

\node[rectangle, draw=awsgray!50, fill=awslightblue, minimum width=\lw,
      minimum height=\lh, font=\small\bfseries, text centered,
      rounded corners=2pt, below=0pt of l2]
  (l1) {Layer 1 — Data Ingestion (ArXiv API + Hacker News API) \quad \textcolor{awsgray}{\scriptsize Batch (weekly) + Near-RT (5 min)}};

\end{tikzpicture}
\caption{Six-Layer Conceptual Architecture of TechPulse. Layers run from data ingestion (bottom) to user interaction and monitoring (top), each with well-defined responsibilities.}
\label{fig:six_layers}
\end{figure}

\section{Alignment with the AWS Well-Architected Framework}
\label{sec:well_architected}

The TechPulse architecture is designed in explicit alignment with the AWS Well-Architected Framework (AWS, 2023), which defines six pillars of cloud system best practice. Table~\ref{tab:well_arch} summarises how each pillar is addressed within the TechPulse design, and Figure~\ref{fig:well_arch} presents a visual mapping of the framework pillars to specific system components.

\begin{table}[H]
\caption{TechPulse Alignment with the AWS Well-Architected Framework Six Pillars.}
\label{tab:well_arch}
\centering
\begin{tabular}{p{3.2cm}p{9.3cm}}
\toprule
\textbf{Pillar} & \textbf{TechPulse Implementation} \\
\midrule
\textbf{Operational Excellence} & Automated ingestion via EventBridge; CI/CD pipeline using GitHub Actions and SAM/CDK; CloudWatch logging and health-check Lambda for proactive monitoring \\[4pt]
\textbf{Security} & Controlled API access via Amazon API Gateway; all services managed and isolated within AWS VPC boundaries; no direct model invocation from the frontend \\[4pt]
\textbf{Reliability} & Decoupled ingestion and embedding pipelines via SQS; Dead Letter Queue for failure recovery and replay; bounded retry logic for Bedrock invocations \\[4pt]
\textbf{Performance Efficiency} & Serverless Lambda scaling; lightweight MiniLM embedding model suited to low-memory serverless execution; metadata pre-filtering reduces vector search scope \\[4pt]
\textbf{Cost Optimization} & Serverless-only execution (pay-per-invocation); AWS Budgets alerts with programmatic halt; estimated prototype cost of USD 10--15/month at 900 queries/month \\[4pt]
\textbf{Sustainability} & Event-driven processing avoids idle compute; SQS decoupling prevents over-provisioning; right-sized Lambda memory allocations \\
\bottomrule
\end{tabular}
\end{table}

% ---- TikZ: Well-Architected Framework mapping diagram ----
\begin{figure}[H]
\centering
\begin{tikzpicture}[node distance=0.35cm and 0.4cm]

% Central TechPulse node
\node[rectangle, rounded corners=6pt, draw=awsblue, fill=awslightblue,
      font=\small\bfseries, minimum height=1.0cm, minimum width=2.8cm,
      text centered] (center) {TechPulse AWS Arch.};

% Six pillar nodes — kept narrower to stay within margins
% Top
\node[rectangle, rounded corners=4pt, draw=awsorange, fill=awslightorange,
      font=\scriptsize\bfseries, minimum height=0.85cm, minimum width=3.4cm,
      text width=3.2cm, text centered,
      above=1.3cm of center] (opex)
  {\textbf{Operational Excellence}\\CI/CD, CloudWatch, EventBridge};

% Bottom
\node[rectangle, rounded corners=4pt, draw=awsgreen, fill=awslightgreen,
      font=\scriptsize\bfseries, minimum height=0.85cm, minimum width=3.4cm,
      text width=3.2cm, text centered,
      below=1.3cm of center] (sust)
  {\textbf{Sustainability}\\Event-driven, right-sized Lambda};

% Left top
\node[rectangle, rounded corners=4pt, draw=propred, fill=propred!10,
      font=\scriptsize\bfseries, minimum height=0.85cm, minimum width=3.2cm,
      text width=3.0cm, text centered,
      left=1.6cm of center, yshift=0.9cm] (sec)
  {\textbf{Security}\\API GW, managed AWS services};

% Left bottom
\node[rectangle, rounded corners=4pt, draw=awspurple, fill=awslightpurple,
      font=\scriptsize\bfseries, minimum height=0.85cm, minimum width=3.2cm,
      text width=3.0cm, text centered,
      left=1.6cm of center, yshift=-0.9cm] (rel)
  {\textbf{Reliability}\\SQS, DLQ, retry logic};

% Right top
\node[rectangle, rounded corners=4pt, draw=awsorange, fill=awslightorange,
      font=\scriptsize\bfseries, minimum height=0.85cm, minimum width=3.2cm,
      text width=3.0cm, text centered,
      right=1.6cm of center, yshift=0.9cm] (perf)
  {\textbf{Performance Efficiency}\\Serverless Lambda, MiniLM};

% Right bottom
\node[rectangle, rounded corners=4pt, draw=propgreen, fill=propgreen!10,
      font=\scriptsize\bfseries, minimum height=0.85cm, minimum width=3.2cm,
      text width=3.0cm, text centered,
      right=1.6cm of center, yshift=-0.9cm] (cost)
  {\textbf{Cost Optimization}\\Pay-per-use, Budget alerts};

% Arrows from center to pillars
\draw[-{Stealth}, thick, color=awsblue!60] (center.north) -- (opex.south);
\draw[-{Stealth}, thick, color=awsblue!60] (center.south) -- (sust.north);
\draw[-{Stealth}, thick, color=awsblue!60] (center.west) -- +(-0.2,0) |- (sec.east);
\draw[-{Stealth}, thick, color=awsblue!60] (center.west) -- +(-0.2,0) |- (rel.east);
\draw[-{Stealth}, thick, color=awsblue!60] (center.east) -- +(0.2,0) |- (perf.west);
\draw[-{Stealth}, thick, color=awsblue!60] (center.east) -- +(0.2,0) |- (cost.west);

\end{tikzpicture}
\caption{TechPulse alignment with the AWS Well-Architected Framework six pillars: Operational Excellence, Security, Reliability, Performance Efficiency, Cost Optimization, and Sustainability.}
\label{fig:well_arch}
\end{figure}

\section{Data Sources}
\label{sec:data_sources}

Two complementary data sources are proposed to balance scholarly authority with real-time relevance.

\begin{table}[H]
\caption{Proposed Data Sources and Roles.}
\label{tab:data_sources}
\centering
\begin{tabular}{p{3.0cm}p{2.5cm}p{4.5cm}p{2.5cm}}
\toprule
\textbf{Source} & \textbf{Type} & \textbf{Role} & \textbf{Frequency} \\
\midrule
ArXiv API & Batch & Scholarly grounding corpus & Weekly \\
Hacker News API & Near-real-time & Industry signal corpus & Every 5 minutes \\
\bottomrule
\end{tabular}
\end{table}

The combination provides complementary temporal and authority profiles: ArXiv contributes rigorously authored research with stable identifiers suitable for deterministic hashing, while Hacker News contributes practitioner signals and community-curated trending discussions. Future extensions may include StackExchange, Wired AI RSS, TechCrunch AI RSS, and the GitHub REST API, though these are deferred to maintain scope and evaluation rigor.

\section{Data Engineering Pipeline}
\label{sec:pipeline}

Figure~\ref{fig:pipeline} presents the proposed end-to-end data engineering pipeline as a schematic flow diagram, complementing the AWS service-level view in Figure~\ref{fig:aws_arch}. The pipeline is designed around four principles: determinism (identical input produces identical output), idempotency (re-running does not create duplicate records), observability (all state transitions are logged), and failure resilience (failures are isolated and recoverable).

% ---- TikZ: Full data engineering pipeline ----
\begin{figure}[H]
\centering
\resizebox{\textwidth}{!}{%
\begin{tikzpicture}[node distance=0.4cm and 0.5cm, font=\small]

% ---- Sources ----
\node[compnode=awsblue] (arxiv) {ArXiv API\\(weekly)};
\node[compnode=awsblue, below=0.35cm of arxiv] (hn) {Hacker News\\(5 min)};

% ---- EventBridge ----
\node[compnode=awsorange, right=0.8cm of arxiv, yshift=-0.55cm] (eb) {EventBridge\\Scheduler};

% ---- Ingestion Lambda ----
\node[compnode=awsorange, right=0.7cm of eb] (inlambda) {Ingestion\\Lambda};

% ---- S3 Raw ----
\node[compnode=awsgreen, above right=0.1cm and 0.7cm of inlambda] (s3raw) {S3 Raw\\(arxiv / hn)};

% ---- SQS ----
\node[compnode=awsgray, right=0.7cm of inlambda, yshift=-0.55cm] (sqs) {SQS\\Queue};

% ---- DLQ ----
\node[compnode=propred!80, below=0.35cm of sqs] (dlq) {Dead Letter\\Queue};

% ---- Preprocess + Embed Lambda ----
\node[compnode=awsorange, right=0.7cm of sqs] (pelambda) {Preprocess \&\\Embed Lambda\\(MiniLM)};

% ---- S3 Processed ----
\node[compnode=awsgreen, above right=0.1cm and 0.7cm of pelambda] (s3proc) {S3 Processed\\(/processed)};

% ---- Aurora ----
\node[compnode=awspurple, right=0.7cm of pelambda, yshift=-0.55cm] (aurora) {Aurora PG\\+ pgvector};

% ---- Retrieval / Bedrock ----
\node[compnode=awsorange, right=0.7cm of aurora] (rag) {RAG\\Orchestrator};
\node[compnode=awsblue, right=0.7cm of rag] (bedrock) {Amazon\\Bedrock};

% ---- API GW + React ----
\node[compnode=awsgreen, right=0.7cm of bedrock] (apigw) {API\\Gateway};
\node[compnode=awsblue, right=0.5cm of apigw] (react) {React\\Frontend};

% ---- Observability ----
\node[compnode=propgray, below=0.8cm of rag, xshift=0.5cm] (cw) {CloudWatch\\Metrics/Logs};
\node[compnode=propgray, right=0.6cm of cw] (budgets) {AWS Budgets\\+ SNS};
\node[compnode=propgray, left=0.6cm of cw] (hc) {Health Check\\Lambda};

% Arrows: sources → eb → inlambda
\draw[parrow] (arxiv.east) -- ++(0.25,0) |- (eb.west);
\draw[parrow] (hn.east)    -- ++(0.25,0) |- (eb.west);
\draw[parrow] (eb) -- (inlambda);
\draw[parrow] (inlambda.east) -- ++(0.25,0) |- (s3raw.west);
\draw[parrow] (inlambda) -- (sqs);
\draw[parrow, dashed, color=propred] (sqs.south) -- (dlq.north)
  node[midway, right, font=\scriptsize, color=propred] {failures};
\draw[parrow] (sqs) -- (pelambda);
\draw[parrow] (pelambda.east) -- ++(0.25,0) |- (s3proc.west);
\draw[parrow] (pelambda) -- (aurora);
\draw[parrow] (aurora) -- (rag);
\draw[parrow] (rag) -- (bedrock);
\draw[parrow] (bedrock) -- (apigw);
\draw[parrow] (apigw) -- (react);

% Observability arrows
\draw[parrow, dashed, color=propgray] (rag.south) -- (cw.north);
\draw[parrow, dashed, color=propgray] (cw.east) -- (budgets.west);
\draw[parrow, dashed, color=propgray] (hc.east) -- (cw.west);

% Layer labels
\node[below=0.5cm of arxiv, font=\scriptsize\bfseries, color=awsblue, xshift=0.4cm]
  {Sources};
\node[below=0.5cm of eb, font=\scriptsize\bfseries, color=awsorange]
  {Trigger};
\node[below=1.2cm of sqs, font=\scriptsize\bfseries, color=awsorange, xshift=0.8cm]
  {Processing};
\node[below=0.5cm of aurora, font=\scriptsize\bfseries, color=awspurple]
  {Storage};
\node[below=0.5cm of rag, font=\scriptsize\bfseries, color=awsorange, xshift=0.3cm]
  {Serving};
\node[below=0.3cm of cw, font=\scriptsize\bfseries, color=propgray]
  {Observability};

\end{tikzpicture}%
}
\caption{Proposed End-to-End Data Engineering and Serving Pipeline for TechPulse. Solid arrows indicate data and control flow; dashed arrows indicate monitoring and observability signals. The pipeline is designed to be deterministic, idempotent, and failure-resilient.}
\label{fig:pipeline}
\end{figure}

\subsection{Ingestion and Decoupling}

Ingestion will be implemented through AWS Lambda functions orchestrated by Amazon EventBridge, enabling both scheduled and event-driven execution. Raw artifacts will be persisted in Amazon S3 under a structured prefix hierarchy (\texttt{/raw/arxiv/}, \texttt{/raw/hn/}, \texttt{/processed/}) to preserve transformation lineage. The ingestion pipeline will be decoupled from the embedding pipeline via Amazon SQS, providing backpressure control, failure isolation, and independent horizontal scaling. Messages that fail processing after retry attempts will be routed to a Dead Letter Queue (DLQ) for controlled replay.

\subsection{Document State Lifecycle}

Each document will progress through a four-state deterministic lifecycle to guarantee idempotency and traceability. Table~\ref{tab:doc_states} defines these states.

\begin{table}[H]
\caption{Document State Transition Model.}
\label{tab:doc_states}
\centering
\begin{tabular}{p{2.0cm}p{3.8cm}p{3.0cm}p{3.0cm}}
\toprule
\textbf{State} & \textbf{Description} & \textbf{Storage} & \textbf{Trigger} \\
\midrule
RAW & Document ingested from API & Amazon S3 & Successful API fetch \\
PROCESSED & Cleaned, normalized, chunked & S3 (\texttt{/processed}) & Preprocessing complete \\
EMBEDDED & Vector representation generated & Aurora (pgvector) & Embedding computed \\
INDEXED & Available for retrieval & Aurora metadata table & Vector committed \\
\bottomrule
\end{tabular}
\end{table}

To guarantee idempotency, canonicalized content fields will be hashed using SHA-256 before processing. For ArXiv documents, the hash covers title, abstract, and publication date; for Hacker News, title, URL, and timestamp. Duplicate hashes will be skipped at ingestion, ensuring consistent corpus integrity across experimental runs.

Figure~\ref{fig:state_machine} illustrates the proposed document state machine.

% ---- TikZ: Document state machine ----
\begin{figure}[H]
\centering
\resizebox{0.96\textwidth}{!}{%
\begin{tikzpicture}[node distance=1.6cm and 2.2cm]

\node[rectangle, rounded corners=5pt, draw=awsblue!70, fill=awslightblue,
      font=\small\bfseries, minimum height=0.7cm, minimum width=2.2cm,
      text centered] (raw) {RAW};

\node[rectangle, rounded corners=5pt, draw=awsgreen!70, fill=awslightgreen,
      font=\small\bfseries, minimum height=0.7cm, minimum width=2.2cm,
      text centered, right=of raw] (proc) {PROCESSED};

\node[rectangle, rounded corners=5pt, draw=awspurple!70, fill=awslightpurple,
      font=\small\bfseries, minimum height=0.7cm, minimum width=2.2cm,
      text centered, right=of proc] (emb) {EMBEDDED};

\node[rectangle, rounded corners=5pt, draw=awsorange!70, fill=awslightorange,
      font=\small\bfseries, minimum height=0.7cm, minimum width=2.2cm,
      text centered, right=of emb] (idx) {INDEXED};

% failure nodes
\node[rectangle, rounded corners=4pt, draw=propred!70, fill=propred!10,
      font=\scriptsize, minimum height=0.55cm, minimum width=1.5cm,
      text centered, below=1.0cm of proc] (dlqbox) {DLQ (replay)};

\draw[parrow] (raw) -- node[above, font=\scriptsize]{Normalize + chunk} (proc);
\draw[parrow] (proc) -- node[above, font=\scriptsize]{MiniLM embed} (emb);
\draw[parrow] (emb) -- node[above, font=\scriptsize]{Commit to Aurora} (idx);

% failure arrows
\draw[-{Stealth}, dashed, color=propred, thick] (proc.south) --
  node[right, font=\scriptsize, color=propred]{fail} (dlqbox.north);
\draw[-{Stealth}, dashed, color=propred, thick] (emb.south) --
  ++(0,-0.5) -| (dlqbox.east);

% replay arrow
\draw[-{Stealth}, dashed, color=propgreen, thick] (dlqbox.west) --
  ++(-0.4,0) |- (raw.south)
  node[near end, below, font=\scriptsize, color=propgreen]{replay};

% SHA note
\node[below=0.15cm of raw, font=\scriptsize, color=awsgray]
  {SHA-256 hash};

\end{tikzpicture}%
}
\caption{Document state machine: RAW $\to$ PROCESSED $\to$ EMBEDDED $\to$ INDEXED. Failures are routed to the Dead Letter Queue for replay, ensuring idempotency.}
\label{fig:state_machine}
\end{figure}

\subsection{Preprocessing and Chunking}

Textual preprocessing will include HTML stripping, encoding normalization, whitespace consolidation, and removal of low-information tokens. Following normalization, documents will be segmented using token-based chunking. An initial window of 500 tokens with 50-token overlap is proposed. This configuration is intended to balance semantic coherence with retrieval granularity: smaller segments may fragment contextual meaning and reduce embedding stability, whereas larger segments risk diluting topical specificity and increasing token consumption during generation. This parameter will be validated empirically during the implementation phase.

\section{Embedding and Indexing}
\label{sec:embedding}

Semantic representations will be generated using the MiniLM sentence-transformer architecture (Wang et al., 2020), selected for its favorable balance of semantic expressiveness and computational efficiency under serverless execution constraints. Each chunk will be converted to a fixed-dimensional embedding and stored in Aurora PostgreSQL using the pgvector extension. Vector similarity search will be implemented using exact cosine similarity to ensure deterministic experimental behavior.

\section{Storage Architecture (Medallion Layout)}
\label{sec:storage}

\begin{table}[H]
\caption{Storage Layers and Responsibilities (Medallion Architecture).}
\label{tab:storage_layers}
\centering
\begin{tabular}{p{2.8cm}p{5.8cm}p{4.0cm}}
\toprule
\textbf{Layer} & \textbf{Technology} & \textbf{Purpose} \\
\midrule
Raw (Bronze) & Amazon S3 (\texttt{/raw/arxiv/}, \texttt{/raw/hn/}) & Store unprocessed API responses \\
Processed (Silver) & Amazon S3 (\texttt{/processed/}) & Store cleaned and chunked documents \\
Vector (Gold) & Aurora PostgreSQL + pgvector & Store embeddings and structured metadata \\
\bottomrule
\end{tabular}
\end{table}

\section{Prompt Construction and Grounded Generation}
\label{sec:prompt_ch4}

Retrieved top-$k$ documents will be concatenated into a structured prompt template that instructs the generative model to produce responses exclusively grounded in the supplied context. The prompt explicitly prohibits unsupported speculation and requires the model to acknowledge insufficient evidence when appropriate. Context length will be truncated to remain within model token limits, ensuring deterministic inference behavior. This approach mitigates hallucination risk and reinforces citation traceability. The generation prompt is organized into four logical sections: (1) \textbf{System Instruction}—defines assistant behavior and enforces evidence-grounded responses; (2) \textbf{Context Block}—concatenates retrieved document excerpts with source identifiers; (3) \textbf{User Question}—the original user query; and (4) \textbf{Citation Instruction}—requires the model to reference supporting documents.

\begin{lstlisting}[caption={Proposed Structured Generation Prompt Template.}, label={lst:prompt}]
SYSTEM:
Answer using only the provided context.
If insufficient evidence exists, state that the
answer cannot be determined.

CONTEXT:
[Document 1 excerpt]
Source: ArXiv

[Document 2 excerpt]
Source: Hacker News

QUESTION:
{User Query}

INSTRUCTION:
Provide a concise answer and cite supporting sources.
\end{lstlisting}

If retrieval returns insufficient evidence, the system will return an abstention response rather than generating unsupported content. The weight parameters and context truncation strategy will be finalized during implementation.

\section{Deployment, Monitoring, and MLOps}
\label{sec:mlops}

Figure~\ref{fig:mlops_loop} presents the proposed MLOps lifecycle loop encompassing continuous ingestion, versioned deployment, observability, and feedback-driven iteration.

% ---- TikZ: MLOps lifecycle loop ----
\begin{figure}[H]
\centering
\resizebox{0.92\textwidth}{!}{%
\begin{tikzpicture}[node distance=1.4cm and 0.5cm]

% Main cycle nodes
\node[pbox=awsblue, minimum width=2.8cm] (ingest) {Data Ingestion\\(EventBridge)};
\node[pbox=awsgreen, right=2.0cm of ingest] (process) {Preprocess \&\\Embed};
\node[pbox=awspurple, right=2.0cm of process] (index) {Index\\(pgvector)};
\node[pbox=awsorange, below=1.8cm of index] (serve) {Serve\\(Bedrock + RAG)};
\node[pbox=propblue, below=1.8cm of ingest] (monitor) {Monitor\\(CloudWatch)};
\node[pbox=propgreen, below=1.8cm of process] (evaluate) {Evaluate\\(RAGAS + Manual)};

% Cycle arrows
\draw[parrow, thick] (ingest) -- (process);
\draw[parrow, thick] (process) -- (index);
\draw[parrow, thick] (index) -- (serve);
\draw[parrow, thick] (serve) -- (evaluate);
\draw[parrow, thick] (evaluate) -- (monitor);
\draw[parrow, thick] (monitor) -- (ingest);

% CI/CD overlay — above the top row
\node[pbox=propgray, above=1.1cm of process, minimum width=5.0cm] (cicd)
  {CI/CD: GitHub Actions + SAM/CDK};
\draw[-{Stealth}, dashed, color=propgray, thick] (cicd.south) -- (process.north);
\draw[-{Stealth}, dashed, color=propgray, thick] (cicd.south) -- ++(0,-0.3) -| (index.north);

% Budget guard — moved BELOW the serve node to avoid right overflow
\node[pbox=propred, below=0.8cm of serve, minimum width=2.6cm] (budget)
  {Budget Guard\\(AWS Budgets)};
\draw[-{Stealth}, dashed, color=propred, thick] (serve.south) -- (budget.north);
\draw[-{Stealth}, dashed, color=propred, thick] (budget.west) -- ++(-0.4,0) |- (monitor.east);

\end{tikzpicture}%
}
\caption{Proposed MLOps lifecycle loop: continuous ingestion, preprocessing, indexing, serving, evaluation, and monitoring, with CI/CD automation and a budget governance guard.}
\label{fig:mlops_loop}
\end{figure}

\subsection{Bedrock Invocation Configuration and Risk Control}

Generation will be performed through Amazon Bedrock using a managed Claude-family model (e.g., Claude 3 Haiku), selected for its low-latency inference profile suitable for real-time retrieval-augmented workflows. To control operational risk and cost exposure, generation parameters will be explicitly bounded. Table~\ref{tab:bedrock_params} presents the proposed parameter configuration.

\begin{table}[H]
\caption{Proposed Bedrock Generation Parameter Configuration.}
\label{tab:bedrock_params}
\centering
\begin{tabular}{lll}
\toprule
\textbf{Parameter} & \textbf{Proposed Value} & \textbf{Rationale} \\
\midrule
Max Output Tokens & 300 & Cost control \\
Temperature & 0.2 & Reduce hallucination risk \\
Timeout & 8 seconds & Prevent latency spikes \\
Retry Attempts & 2 & Transient failure handling \\
\bottomrule
\end{tabular}
\end{table}

Retry logic will follow a bounded exponential backoff policy. If invocation fails after two retry attempts, the system will degrade gracefully by returning a retrieval-only summary constructed from the top-$k$ retrieved documents without generative expansion, and the error will be logged to CloudWatch. Bedrock invocation will be configured within the Asia-Pacific region (\texttt{ap-southeast-1}) to minimize latency. In the event of temporary regional unavailability, the system will return a grounded retrieval summary rather than attempting cross-region invocation, preserving latency constraints and compliance considerations.

AWS Budgets alerts will programmatically halt invocations if predefined monthly cost thresholds are approached. This budget guard mechanism prevents uncontrolled token consumption and ensures predictable operational expenditure.

\subsection{Bedrock Backup Plan and Fallback Strategy}
\label{sec:bedrock_backup}

Amazon Bedrock is the primary inference service for TechPulse. However, Bedrock is not available in the Learner Lab environment, and regional availability may be subject to access restrictions. To ensure the system remains functional under these constraints, the following fallback strategy is defined.

\textbf{Primary option (Bedrock — Claude Haiku).} Preferred for low-latency managed inference within the \texttt{ap-southeast-1} region. This path will be used for final evaluation if Bedrock access is confirmed.

\textbf{Fallback Option 1 — Ollama with a locally hosted model.} If Bedrock access is unavailable, generation will be redirected to a locally hosted model via Ollama (e.g., Llama 3.1 8B or Mistral 7B). Ollama provides an OpenAI-compatible REST API, so the RAG orchestration Lambda can be reconfigured by changing a single endpoint environment variable with no changes to the retrieval pipeline or evaluation logic. The RAG orchestrator will be implemented with a configurable \texttt{LLM\_BACKEND} environment variable (\texttt{bedrock} or \texttt{ollama}) to enable transparent switching.

\textbf{Fallback Option 2 — Hugging Face Inference API.} If neither Bedrock nor local inference is feasible, the Hugging Face Inference API provides hosted access to open-weight models (e.g., Mistral-7B-Instruct) under a free tier suitable for prototype-scale evaluation. This requires no local compute and incurs minimal cost within evaluation query volumes.

All three options produce responses compatible with the same structured prompt template (Listing~\ref{lst:prompt}), ensuring that evaluation results remain comparable across backends. Table~\ref{tab:llm_backends} summarizes the fallback hierarchy.

\begin{table}[H]
\caption{LLM Backend Fallback Hierarchy.}
\label{tab:llm_backends}
\centering
\begin{tabular}{p{0.5cm}p{3.5cm}p{4.0cm}p{4.0cm}}
\toprule
\textbf{Tier} & \textbf{Backend} & \textbf{Condition} & \textbf{Notes} \\
\midrule
1 & Amazon Bedrock (Claude Haiku) & Bedrock access confirmed & Primary path; managed, low-latency \\
2 & Ollama (local) & Bedrock unavailable & OpenAI-compatible API; config-only switch \\
3 & HuggingFace Inference API & Local compute unavailable & Free tier; suitable for eval-scale queries \\
\bottomrule
\end{tabular}
\end{table}

\subsection{Data Drift Detection and Remediation Actions}
\label{sec:drift}

Because TechPulse ingests data continuously from a rapidly evolving technology domain, the statistical properties of the indexed corpus may shift over time—a phenomenon known as data drift. Left unaddressed, drift can cause context precision to degrade even if the retrieval pipeline itself remains unchanged. TechPulse incorporates a lightweight drift monitoring strategy and a defined set of remediation actions.

\textbf{Drift detection mechanism.} At each weekly ingestion cycle, a reference distribution of cosine similarity scores will be computed over a fixed set of 20 canonical probe queries against the current indexed corpus. This distribution is persisted in CloudWatch as a time-series metric. If the mean cosine similarity for the probe set drops by more than 10\% relative to the previous week's baseline—or if a statistical process control alert (Shewhart $3\sigma$ rule) is triggered—a drift alert is raised via CloudWatch alarm and routed through SNS to the operational team.

\textbf{Remediation actions on drift detection.} Upon drift alert, the following ordered response protocol will be applied:

\begin{enumerate}[leftmargin=0.5in]
  \item \textbf{Corpus audit.} The ingestion logs will be inspected to identify whether the drift is localized to a specific source (ArXiv or Hacker News) or a specific topic cluster. Localized drift may indicate a change in API response format, a shift in publication volume, or a sudden influx of off-topic documents.

  \item \textbf{Metadata filter re-calibration.} If the drift is associated with off-topic documents entering the corpus, the metadata pre-filter constraints (recency window, keyword whitelist) will be tightened to restore corpus topical coherence without re-indexing.

  \item \textbf{Re-indexing the affected corpus segment.} If drift affects the embedding space (e.g., new terminology that is underrepresented in the current embedding model vocabulary), the affected document segments will be re-chunked and re-embedded with updated preprocessing parameters and re-indexed into Aurora.

  \item \textbf{Embedding model refresh (deferred).} Full embedding model retraining or replacement is deferred to a post-prototype phase, as it requires significant compute and evaluation overhead beyond the current project scope.
\end{enumerate}

This protocol ensures that drift is detected promptly and addressed through targeted, incremental remediation rather than full system redeployment, preserving operational continuity.

A scheduled health-check Lambda will be triggered periodically via EventBridge to validate: (1) Aurora connectivity via a lightweight vector similarity query; (2) Bedrock accessibility via a minimal test prompt; and (3) API Gateway round-trip latency. Results will be logged to CloudWatch, with failure conditions triggering CloudWatch alarms for operational notification. This proactive monitoring strategy enables early detection of dependency degradation and prevents silent system failures. Together with DLQ monitoring and budget guard mechanisms, these controls ensure reliability, fault isolation, and operational resilience.

\subsection{CI/CD and Deployment Automation}

All application code—including ingestion functions, preprocessing logic, retrieval orchestration, and frontend components—will be version-controlled in a GitHub repository. A CI/CD pipeline implemented using GitHub Actions will automatically execute linting and unit tests upon commit, build deployable artifacts, and deploy updates to AWS services. Infrastructure components will be defined using Infrastructure-as-Code principles (AWS SAM or CDK) to ensure reproducible and consistent deployments across environments. This automated pipeline reduces manual deployment risk, enforces code quality checks, and supports controlled iteration of retrieval logic, prompt templates, and orchestration workflows.The project source code will be maintained in a GitHub organization repository to support version control, team collaboration, and deployment automation.

\section{User Interface}
\label{sec:ui}

The user interface will be a lightweight React single-page application hosted on Amazon S3. It will provide natural language query submission, structured response rendering, citation visualization, and retrieval metadata transparency. All responses will be constrained through the RAG pipeline; direct model invocation will not be permitted. When retrieval evidence is insufficient, the system will return an explicit fail-safe message.

% ============================================================
% CHAPTER 5: EXPERIMENTAL PLAN AND EVALUATION METRICS
% ============================================================
\chapter{Experimental Plan and Evaluation Metrics}

\section{Evaluation Objective}
\label{sec:eval_objective}

The evaluation aims to answer three questions: (1) Does hybrid retrieval improve contextual precision and grounding faithfulness compared to vector-only retrieval? (2) What latency and cost overhead does hybrid retrieval introduce? (3) Does the system remain within acceptable operational thresholds under moderate workload?

\section{Experimental Configurations}
\label{sec:eval_configs}

Two configurations will be compared under identical corpus, embedding model, and chunking parameters:

\begin{table}[H]
\caption{Experimental Configuration Comparison.}
\label{tab:exp_configs}
\centering
\begin{tabular}{p{3.5cm}cc}
\toprule
\textbf{Component} & \textbf{Baseline (Vector-Only)} & \textbf{Hybrid (Proposed)} \\
\midrule
Metadata Filtering & $\times$ & \checkmark \\
Vector Similarity & \checkmark & \checkmark \\
Reranking-Lite & $\times$ & \checkmark \\
Recency Weighting & $\times$ & \checkmark \\
Top-$k$ & 5 & 5 \\
Embedding Model & MiniLM & MiniLM \\
Cost Governance & None & AWS Budgets guard \\
\bottomrule
\end{tabular}
\end{table}

\section{Evaluation Dataset}
\label{sec:eval_dataset}

Evaluation will use approximately 20--30 manually curated trend-oriented queries spanning recent research developments, emerging tools, and comparative technology discussions in artificial intelligence and related domains. Example queries include: ``What are recent approaches to efficient transformer architectures?''; ``What emerging AI tools are currently being discussed?''; and ``What are recent trends in large language model compression?''

\textbf{Justification for evaluation set size.} The choice of 20--30 queries is motivated by three considerations. First, the evaluation objective is a \emph{controlled comparison} between two retrieval configurations (baseline vector-only vs.\ hybrid) under identical corpus and embedding conditions—not an open-ended benchmark generalization study. In this setting, the relevant statistical question is whether the hybrid configuration consistently outperforms the baseline, not whether absolute scores generalize across all possible queries. A paired comparison across 20--30 carefully constructed queries is sufficient to detect a meaningful precision difference if one exists, consistent with the evaluation scale adopted in prior RAG evaluation work (Es et al., 2023; Saad-Falcon et al., 2023). Second, the queries will be stratified across three query types—(a) recent research trend queries, (b) specific tool or technology queries, and (c) comparative or synthesis queries—ensuring that results are not driven by a single query category. Third, at the prototype stage, manual annotation quality is more important than annotation volume: a smaller, carefully annotated set with verified relevance judgments produces more reliable precision estimates than a larger set with noisier annotations. Inter-annotator agreement will be computed using Cohen's $\kappa$ across a 10-query overlap sample annotated independently by both team members; a $\kappa > 0.7$ will be required to proceed with the full evaluation.

Retrieved documents will be manually annotated using a two-level relevance scheme: \textbf{Relevant} (directly addresses the query) and \textbf{Partially Relevant} (topically related but not directly answering).

\section{Evaluation Metrics}
\label{sec:eval_metrics}

Figure~\ref{fig:eval_framework} summarizes the proposed evaluation framework.

% ---- TikZ: Evaluation framework diagram ----
\begin{figure}[H]
\centering
\resizebox{\textwidth}{!}{%
\begin{tikzpicture}

% Use matrix layout to ensure proper centering
\node[evalbox=awsblue] (prec) at (-6.0,0)
  {\textbf{Context Precision}\\[2pt]\scriptsize Relevant / Total docs\\Target: $> 0.75$};

\node[evalbox=awsgreen] (faith) at (-2.0,0)
  {\textbf{Faithfulness}\\[2pt]\scriptsize Claims grounded in\\retrieved context\\High target};

\node[evalbox=awsorange] (lat) at (2.0,0)
  {\textbf{p95 Latency}\\[2pt]\scriptsize End-to-end response\\(CloudWatch)\\$< 2$ sec};

\node[evalbox=awspurple] (cost) at (6.0,0)
  {\textbf{Cost / Query}\\[2pt]\scriptsize Bedrock + Lambda\\+ Aurora\\$<$USD 10/mo};

\node[rectangle, draw=awsgray!40, fill=sectionbg, rounded corners=3pt,
      font=\scriptsize, minimum height=0.55cm, text width=3.2cm, text centered]
  (m1) at (-6.0,-1.7) {Manual annotation\\+ RAGAS};

\node[rectangle, draw=awsgray!40, fill=sectionbg, rounded corners=3pt,
      font=\scriptsize, minimum height=0.55cm, text width=3.2cm, text centered]
  (m2) at (-2.0,-1.7) {Citation traceability\\+ hallucination check};

\node[rectangle, draw=awsgray!40, fill=sectionbg, rounded corners=3pt,
      font=\scriptsize, minimum height=0.55cm, text width=3.2cm, text centered]
  (m3) at (2.0,-1.7) {AWS CloudWatch\\Logs};

\node[rectangle, draw=awsgray!40, fill=sectionbg, rounded corners=3pt,
      font=\scriptsize, minimum height=0.55cm, text width=3.2cm, text centered]
  (m4) at (6.0,-1.7) {AWS Billing /\\Usage Metrics};

\draw[parrow] (prec) -- (m1);
\draw[parrow] (faith) -- (m2);
\draw[parrow] (lat) -- (m3);
\draw[parrow] (cost) -- (m4);

\end{tikzpicture}%
}
\caption{Proposed evaluation framework: four dimensions—context precision, faithfulness, p95 latency, and cost per query—each with a measurement method and target threshold.}
\label{fig:eval_framework}
\end{figure}

\subsection{Context Precision}

Retrieval precision will be measured as:

\begin{equation}
  \text{Precision} = \frac{\text{Number of relevant retrieved documents}}{\text{Total retrieved documents}}
  \label{eq:precision}
\end{equation}

The preliminary target threshold is $>0.75$, subject to revision based on empirical findings.

\subsection{Faithfulness and Hallucination Verification}

Faithfulness measures whether generated responses remain grounded in retrieved context. As detailed in Section~\ref{sec:prompt}, TechPulse applies a three-layer hallucination verification approach: (1) prompt-level constraints enforcing evidence grounding and abstention; (2) RAGAS faithfulness scoring decomposing responses into atomic claims and checking each against retrieved documents, with a threshold of 0.8 triggering manual review; and (3) citation traceability enforcement intercepting unanchored claims at the output layer. Assessment will therefore combine automated RAGAS faithfulness scoring across all evaluation queries and manual inspection of flagged responses where RAGAS faithfulness falls below threshold.

\subsection{Latency}

End-to-end p95 latency will be measured from API request submission to response return, decomposed into query embedding time, retrieval time, reranking time, and Bedrock generation time. The preliminary threshold is p95 $< 2$ seconds.

\subsection{Cost Per Query and Projected Monthly Cost}

Cost per query will be estimated based on: Bedrock token usage (input and output), Lambda execution duration, and Aurora query execution cost. To ensure operational feasibility, a preliminary cost estimation is provided based on expected moderate system usage.

Assuming a moderate workload of approximately 10 users per day with three queries per user, the system is expected to process approximately 900 queries per month. Under this workload, serverless compute costs remain minimal due to short Lambda execution durations, while Bedrock token usage constitutes the primary variable cost. Based on preliminary estimates using current AWS pricing, the overall monthly operational cost is expected to remain in a low-cost research prototype range of approximately USD~10--15 per month. Actual costs may vary depending on usage patterns, context length, and model token consumption. Detailed cost measurement will be performed using AWS billing metrics during the evaluation phase.

The budget guard mechanism described in Section~\ref{sec:mlops} will programmatically halt Bedrock invocations if the monthly threshold is approached, preventing cost overruns during the evaluation period.

\begin{table}[H]
\caption{Summary of Evaluation Metrics and Preliminary Targets.}
\label{tab:eval_targets}
\centering
\begin{tabular}{p{3.0cm}p{3.0cm}p{3.2cm}p{3.0cm}}
\toprule
\textbf{Metric} & \textbf{Dimension} & \textbf{Method} & \textbf{Preliminary Target} \\
\midrule
Context Precision & Retrieval quality & Manual + RAGAS & $> 0.75$ \\
Faithfulness & Grounding quality & Citation traceability & High (qualitative) \\
p95 Latency & System performance & CloudWatch logs & $< 2$ seconds \\
Cost per Query & Operational cost & AWS billing metrics & $<$~USD 15/month \\
\bottomrule
\end{tabular}
\end{table}

\textit{Note: All target values are preliminary and will be revisited based on initial empirical results. The evaluation methodology may be refined as the project progresses.}

\section{Scalability Analysis}
\label{sec:scalability}

Scalability will be assessed analytically rather than through large-scale stress testing at the prototype stage. Vector similarity search in pgvector operates at $O(Nd)$ complexity. As corpus size grows, ANN indexing strategies and Aurora Serverless auto-scaling may be evaluated as extensions.

\section{Discussion Framework}
\label{sec:discussion}

The analysis will interpret: precision improvements versus computational overhead; trade-offs between retrieval strictness and recall; cost--performance balance; and suitability for production deployment. The emphasis will be placed on demonstrating that lightweight hybrid retrieval can improve grounding quality without violating latency or cost constraints.

% ============================================================
% CHAPTER 6: EXPECTED CONTRIBUTIONS
% ============================================================
\chapter{Expected Contributions}

\section{Summary of Contributions}
\label{sec:contributions}

This project is expected to make the following contributions, subject to the empirical findings that will emerge during implementation and evaluation.

\textbf{Technical Contributions.}

\begin{enumerate}[leftmargin=0.5in]
  \item \textit{Hybrid RAG architecture with recency-aware reranking.} A lightweight, interpretable reranking function combining semantic similarity, lexical overlap, and exponential recency decay, designed specifically for dynamic knowledge domains where temporal relevance is critical.

  \item \textit{Production-grade MLOps integration.} A demonstration that RAG systems can be deployed with full operational rigor—deterministic ingestion pipelines, idempotent state tracking, CI/CD automation, health monitoring, and cost governance—within a serverless AWS environment, bridging the gap between research prototypes and production systems.

  \item \textit{Unified vector--relational storage without dedicated infrastructure.} Evidence that Aurora PostgreSQL with pgvector provides sufficient performance for hybrid retrieval at prototype scale, avoiding the operational complexity of dedicated vector databases.

  \item \textit{Empirical evaluation of hybrid vs.\ baseline retrieval.} A controlled comparison under identical corpus and embedding conditions, quantifying the precision and grounding gains from metadata filtering and reranking-lite against the associated latency and cost overhead.
\end{enumerate}

\textbf{Practical Contributions.}

\begin{enumerate}[leftmargin=0.5in]
  \item A reusable architecture blueprint for production-grade, cost-governed hybrid RAG systems deployable on serverless AWS infrastructure.

  \item An evaluation methodology for assessing RAG systems across retrieval quality, grounding faithfulness, operational latency, and cost dimensions simultaneously.

  \item A practical demonstration that emerging technology intelligence synthesis can be grounded, citation-transparent, and operationally maintainable within a research project budget.
\end{enumerate}

Figure~\ref{fig:contributions} summarizes the positioning of TechPulse across the key contribution dimensions.

% ---- TikZ: Contribution positioning summary ----
\begin{figure}[H]
\centering
\begin{tikzpicture}[node distance=0.3cm]

\def\bw{12.0cm}

\node[rectangle, draw=awsblue!60, fill=awslightblue, minimum width=\bw,
      minimum height=0.7cm, font=\small, text centered, rounded corners=3pt]
  (c1) {Hybrid Retrieval with Recency-Aware Reranking \quad \textcolor{awsblue}{\scriptsize NLP / Retrieval}};

\node[rectangle, draw=awsgreen!60, fill=awslightgreen, minimum width=\bw,
      minimum height=0.7cm, font=\small, text centered, rounded corners=3pt,
      below=0.2cm of c1]
  (c2) {Production-Grade MLOps Integration (AWS-native) \quad \textcolor{awsgreen}{\scriptsize MLOps}};

\node[rectangle, draw=awspurple!60, fill=awslightpurple, minimum width=\bw,
      minimum height=0.7cm, font=\small, text centered, rounded corners=3pt,
      below=0.2cm of c2]
  (c3) {Unified Vector--Relational Storage (pgvector) \quad \textcolor{awspurple}{\scriptsize Data Engineering}};

\node[rectangle, draw=awsorange!60, fill=awslightorange, minimum width=\bw,
      minimum height=0.7cm, font=\small, text centered, rounded corners=3pt,
      below=0.2cm of c3]
  (c4) {Empirical Hybrid vs.\ Baseline Retrieval Comparison \quad \textcolor{awsorange}{\scriptsize Empirical Research}};

\node[rectangle, draw=propblue!60, fill=propblue!10, minimum width=\bw,
      minimum height=0.7cm, font=\small, text centered, rounded corners=3pt,
      below=0.2cm of c4]
  (c5) {Multi-Dimensional Evaluation: Precision + Faithfulness + Latency + Cost \quad \textcolor{propblue}{\scriptsize Evaluation}};

\end{tikzpicture}
\caption{Summary of expected contributions across five dimensions, spanning NLP/retrieval, MLOps, data engineering, empirical research, and evaluation methodology.}
\label{fig:contributions}
\end{figure}

\section{Source Code Availability}
\label{sec:code_availability}

To promote reproducibility and transparent system development, the source code for the proposed TechPulse platform will be maintained in a GitHub organization repository. The repository is available at:

\begin{center}
\url{https://github.com/orgs/EZ2-555555/repositories}
\end{center}

This repository is intended to host the major implementation artifacts of the project, including data ingestion modules, preprocessing and chunking logic, embedding and retrieval components, RAG orchestration code, infrastructure-as-code definitions, and the user interface frontend. It will also support collaborative development and deployment automation through GitHub Actions as part of the proposed CI/CD workflow.

Because this report represents the proposal stage, the repository contents and internal structure may be refined during implementation. The final project report will provide the definitive repository structure, documentation, and implementation details.


\section{Limitations and Scope Boundaries}
\label{sec:limitations}

This proposal acknowledges the following scope boundaries: (1) the evaluation corpus is limited to ArXiv and Hacker News at prototype scale; generalizability to other domains or larger corpora requires further study; (2) the evaluation query set of 20–30 curated queries is sufficient for preliminary comparative analysis, consistent with prior RAG and dense-retrieval evaluation practices where compact but carefully annotated query sets can produce statistically meaningful comparisons when experimental variables are controlled (Es et al., 2023). However, such a query set may not fully capture the diversity of real-world user information needs. If the system were extended beyond the prototype stage, the evaluation dataset could be expanded to include a larger and more diverse set of queries drawn from additional technology domains and user scenarios; (3) cross-encoder reranking and ANN indexing are out of scope for this project but identified as future research directions; and (4) specific implementation details—chunking parameters, reranking weights, temporal window thresholds—will be tuned empirically and may deviate from the initial values proposed here.

% ============================================================
% REFERENCES
% ============================================================
\chapter*{References}
\addcontentsline{toc}{chapter}{REFERENCES}

\begin{flushleft}
\setlength{\parskip}{8pt}
\setlength{\hangindent}{0.5in}

AWS. (2023). \textit{AWS well-architected framework}. Amazon Web Services. \url{https://aws.amazon.com/architecture/well-architected/}

Bayram, F., Ahmed, B. S., \& Kassler, A. (2022). From concept drift to model degradation: An overview on performance-aware drift detectors. \textit{Knowledge-Based Systems}, \textit{245}, 108632. \url{https://doi.org/10.1016/j.knosys.2022.108632}

Borgeaud, S., Mensch, A., Hoffmann, J., Cai, T., Rutherford, E., Millican, K., van den Driessche, G., Lespiau, J.-B., Damoc, B., Clark, A., de Las Casas, D., Guy, A., Menick, J., Ring, R., Hennigan, T., Huang, S., Maggiore, L., Jones, C., Cassirer, A., \ldots Sifre, L. (2022). Improving language models by retrieving from trillions of tokens. In \textit{Proceedings of the 39th International Conference on Machine Learning} (Vol. 162, pp. 2206--2240). PMLR.

Breck, E., Polyzotis, N., Roy, S., Whang, S., \& Zinkevich, M. (2017). The ML test score: A rubric for ML production readiness and technical debt reduction. In \textit{Proceedings of the IEEE International Conference on Big Data} (pp. 1123--1132). IEEE.

Brown, T., Mann, B., Ryder, N., Subbiah, M., Kaplan, J. D., Dhariwal, P., Neelakantan, A., Shyam, P., Sastry, G., Askell, A., Agarwal, S., Herbert-Voss, A., Krueger, G., Henighan, T., Child, R., Ramesh, A., Ziegler, D., Wu, J., Winter, C., \ldots Amodei, D. (2020). Language models are few-shot learners. \textit{Advances in Neural Information Processing Systems}, \textit{33}, 1877--1901.

Es, S., James, J., Anke, L. E., \& Schockaert, S. (2023). RAGAS: Automated evaluation of retrieval augmented generation. \textit{arXiv preprint arXiv:2309.15217}.

Gama, J., \v{Z}liobait\.{e}, I., Bifet, A., Pechenizkiy, M., \& Bouchachia, A. (2014). A survey on concept drift adaptation. \textit{ACM Computing Surveys}, \textit{46}(4), 1--37. \url{https://doi.org/10.1145/2523813}

Gao, Y., Xiong, Y., Gao, X., Jia, K., Pan, J., Bi, Y., Dai, Y., Sun, J., Wang, M., \& Wang, H. (2023). Retrieval-augmented generation for large language models: A survey. \textit{arXiv preprint arXiv:2312.10997}.

Izacard, G., Lewis, P., Lomeli, M., Hosseini, L., Petroni, F., Schick, T., Dwivedi-Yu, J., Joulin, A., Riedel, S., \& Grave, E. (2022). Atlas: Few-shot learning with retrieval augmented language models. \textit{Advances in Neural Information Processing Systems}, \textit{35}, 22251--22265.

Ji, Z., Lee, N., Frieske, R., Yu, T., Su, D., Xu, Y., Ishii, E., Bang, Y. J., Madotto, A., \& Fung, P. (2023). Survey of hallucination in natural language generation. \textit{ACM Computing Surveys}, \textit{55}(12), 1--38. \url{https://doi.org/10.1145/3571730}

Lewis, P., Perez, E., Piktus, A., Petroni, F., Karpukhin, V., Goyal, N., K{\"u}ttler, H., Lewis, M., Yih, W.-T., Rockt{\"a}schel, T., Riedel, S., \& Kiela, D. (2020). Retrieval-augmented generation for knowledge-intensive NLP tasks. \textit{Advances in Neural Information Processing Systems}, \textit{33}, 9459--9474.

Lu, J., Liu, A., Dong, F., Gu, F., Gama, J., \& Zhang, G. (2018). Learning under concept drift: A review. \textit{IEEE Transactions on Knowledge and Data Engineering}, \textit{31}(12), 2346--2363. \url{https://doi.org/10.1109/TKDE.2018.2876857}

Manakul, P., Liusie, A., \& Gales, M. J. F. (2023). SelfCheckGPT: Zero-resource black-box hallucination detection for generative large language models. In \textit{Proceedings of the 2023 Conference on Empirical Methods in Natural Language Processing} (pp. 9004--9017). Association for Computational Linguistics.

Min, S., Krishna, K., Lyu, X., Lewis, M., Yih, W.-T., Koh, P. W., Iyyer, M., Zettlemoyer, L., \& Hajishirzi, H. (2023). FActScoring: Fine-grained atomic evaluation of factual precision in long form text generation. In \textit{Proceedings of the 2023 Conference on Empirical Methods in Natural Language Processing} (pp. 12076--12100). Association for Computational Linguistics.

Nogueira, R., \& Cho, K. (2019). Passage re-ranking with BERT. \textit{arXiv preprint arXiv:1901.04085}.

Saad-Falcon, J., Khattab, O., Potts, C., \& Zaharia, M. (2023). ARES: An automated evaluation framework for retrieval-augmented generation systems. \textit{arXiv preprint arXiv:2311.09476}.

Sculley, D., Holt, G., Golovin, D., Davydov, E., Phillips, T., Ebner, D., Chaudhary, V., Young, M., Crespo, J.-F., \& Dennison, D. (2015). Hidden technical debt in machine learning systems. \textit{Advances in Neural Information Processing Systems}, \textit{28}, 2503--2511.

Shi, F., Chen, X., Misra, K., Scales, N., Dohan, D., Chi, E. H., Sch\"{a}rli, N., \& Zhou, D. (2023). Large language models can be easily distracted by irrelevant context. In \textit{Proceedings of the 40th International Conference on Machine Learning} (pp. 31210--31227). PMLR.

Thakur, N., Reimers, N., Daxenberger, J., \& Gurevych, I. (2021). BEIR: A heterogeneous benchmark for zero-shot evaluation of information retrieval models. In \textit{Proceedings of the 35th Conference on Neural Information Processing Systems---Datasets and Benchmarks Track}.

Wang, W., Wei, F., Dong, L., Bao, H., Yang, N., \& Zhou, M. (2020). MiniLM: Deep self-attention distillation for task-agnostic compression of pre-trained transformers. \textit{Advances in Neural Information Processing Systems}, \textit{33}, 5776--5788.

Zhang, T., Chen, Y., \& Zhou, Z. (2020). Machine learning operations (MLOps): Overview, definition, and architecture. \textit{IEEE Software}, \textit{37}(6), 1--17.

Zhang, Y., Chen, X., Zhang, Y., Hu, X., Liu, Z., \& Tang, J. (2024). A comprehensive survey on retrieval-augmented generation. \textit{ACM Transactions on Information Systems}. Advance online publication. \url{https://doi.org/10.1145/3703155}


% ============================================================
% APPENDIX A: SOURCE CODE AVAILABILITY
% ============================================================

\appendix

\chapter{Source Code Availability}

The implementation of the proposed \textit{TechPulse: Hybrid Retrieval-Augmented Generation System for Emerging Technology Intelligence} will be developed primarily using Python and modern cloud-native technologies. 

To support reproducibility, transparency, and collaborative development, the project source code will be maintained in a GitHub organization repository.

The complete source code repository is available at:

\begin{center}
\url{https://github.com/orgs/EZ2-555555/repositories}
\end{center}

The repository is expected to include the following major components:

\begin{itemize}

\item Data ingestion modules for collecting data from the ArXiv API and Hacker News API

\item Preprocessing and document chunking scripts for text normalization and token-based segmentation

\item Semantic embedding pipeline using the MiniLM sentence-transformer model

\item Hybrid retrieval implementation including metadata filtering, cosine similarity search, and reranking-lite scoring

\item Retrieval-Augmented Generation (RAG) orchestration components integrating the retrieval pipeline with the generative model backend

\item Infrastructure-as-Code definitions for deploying AWS resources using SAM or CDK

\item Frontend interface implementation using a React-based web application

\item Evaluation scripts for measuring context precision, faithfulness, latency, and cost metrics

\end{itemize}

The repository will also support CI/CD automation through GitHub Actions to enable controlled versioning, automated testing, and reproducible deployment workflows.

As this document represents a project proposal, the repository structure and implementation details may evolve during the development phase. The final project report will document the complete repository structure, system implementation, and experimental evaluation scripts.

The source code repository is provided to ensure transparency, reproducibility, and extensibility of the proposed TechPulse system.


\end{flushleft}

\end{document}

